\section{Lefschetz operators and the comparison criterion}
\label{sec:lefschetz}

For the proof of Theorem~\ref{thm:main}, assume that $n\ge4$ is composite.
We first fix the perverse normalization and identify a common candidate for
the two Lefschetz operators.  We then give the linear-algebra argument that
will compare the filtrations once a common lowering operator has been
constructed.  Throughout this section, cohomology has rational coefficients.

\subsection{The normalized perverse filtration}

Recall that $X$ parametrizes stable, trace-free Higgs bundles of rank $n$
and determinant $L$ on $C$, and that $h:X\to A$ is its Hitchin map.
The standard geometry of this system gives a proper surjective morphism
from a smooth quasi-projective variety, together with a holomorphic
symplectic form of weight one for Higgs-field scaling;
see \cite[Section~1.1]{dCMS2022}.

\begin{lemma}\label{norm:dimensions}
Set $R=(n^2-1)(g-1)$.  Then $\dim X=2R$, $\dim A=R$, and every fiber of
$h$ has dimension $R$.  In particular,
\[
 r:=\dim(X\times_A X)-\dim X=R.
\]
\end{lemma}

\begin{proof}
Riemann--Roch gives
\[
 \dim A=\sum_{i=2}^n h^0(C,K_C^i)
 =\sum_{i=2}^n(2i-1)(g-1)=R.
\]
The trace-free Higgs deformation complex has vanishing zeroth and second
hypercohomology, by stability and Serre duality.  Its Euler characteristic
is $-2R$, so the tangent space has dimension $2R$ at every point.

We recall why the fiber bound holds also over the singular Hitchin values.
Higgs-field scaling makes $X$ semiprojective: properness of $h$ extends
each scaling orbit at zero, since its Hitchin value extends there, and
the fixed locus is closed in the proper fiber $h^{-1}(0)$.
For the Bia\l{}ynicki--Birula decomposition \cite{BialynickiBirula1973}
in this semiprojective setting, see \cite[Section~1.2]{HRV2015}.
At a fixed point, the symplectic form pairs a tangent weight $w$ with
weight $1-w$.  The weights are integers, so the nonpositive tangent
weights account for exactly half the tangent dimension.  Consequently
each repelling stratum, including its fixed component, has dimension
$R$.

The union of these strata is $h^{-1}(0)$.  Indeed, a point with a limit
at infinity must have zero Hitchin value, because every Hitchin coordinate
has positive scaling weight.  Conversely, an orbit in $h^{-1}(0)$ has a
limit at infinity because that fiber is proper.  Thus
$\dim h^{-1}(0)=R$.  Pulling $h$ back along the weighted scaling curve
from any $a\in A$ to zero, upper semicontinuity of fiber dimension for a
proper morphism gives $\dim h^{-1}(a)\le R$.  The dimension inequality
for a morphism from a smooth $2R$-dimensional variety to an
$R$-dimensional variety gives the reverse inequality on each nonempty
fiber.  Surjectivity therefore proves the asserted equidimensionality.
It follows that $\dim(X\times_A X)=R+2R=3R$.
\end{proof}

Put
\[
 K=Rh_*\mathbb Q_X[R].
\]
We normalize the perverse filtration by
\begin{equation}\label{norm:perverse-filtration}
 P_kH^m(X)=\operatorname{Im}\left\{
 H^{m-R}\bigl(A,{}^p\tau_{\le k}K\bigr)
 \longrightarrow H^m(X)\right\}.
\end{equation}
This is the normalization of \cite[Section~1.1]{MS2024}.
The fiber bound implies that $Rh_*\mathbb Q_X[2R]$ has perverse
cohomology only in degrees $[-R,R]$: the upper bound follows from the
support criterion, and the lower bound follows by Verdier duality.
Indeed its ordinary stalk cohomology vanishes in positive degrees,
and every support has dimension at most $R$.
After shifting by $[-R]$, the perverse degrees of $K$ lie in $[0,2R]$.
In particular,
\begin{equation}\label{norm:perverse-range}
 P_{-1}H^*(X)=0,\qquad P_{2R}H^*(X)=H^*(X).
\end{equation}

We will use the following shift convention repeatedly.  A morphism
$K\to K[s]$ induces a map of cohomological degree $s$ sending $P_k$ into
$P_{k+s}$, since
\begin{equation}\label{norm:shift-rule}
 {}^p\tau_{\le k}(K[s])
 =({}^p\tau_{\le k+s}K)[s].
\end{equation}
Thus an actual morphism $Rh_*\mathbb Q_X\to Rh_*\mathbb Q_X[-2]$
will supply precisely the lowering bound needed below.

\subsection{The degree-two class}

The universal projective bundle on $C\times X$ determines the characteristic class
\[
 \widetilde{\operatorname{ch}}_2(E)
 =\operatorname{ch}_2(E)-\frac{c_1(E)^2}{2n}
 =\frac{1}{2n}\operatorname{ch}_2(\operatorname{End}E).
\]
The last expression also defines it when a universal vector bundle has
not been chosen.  Define
\begin{equation}\label{norm:e}
 e=\int_C\widetilde{\operatorname{ch}}_2(E)\in H^2(X).
\end{equation}
We also write $e$ for cup product by this class.

\begin{lemma}\label{norm:low-cohomology}
For the composite ranks $n\ge4$ under consideration,
\[
 H^0(X)=\mathbb Q,\qquad H^1(X)=0,\qquad H^2(X)=\mathbb Q e,
 \qquad e\ne0.
\]
Every ample class on $X$ is a nonzero scalar multiple of $e$.
\end{lemma}

\begin{proof}
The finite group $\Gamma=\operatorname{Pic}^0(C)[n]$ acts on $X$ by
tensor product.  Averaging over $\Gamma$ decomposes rational cohomology
into its invariant summand and its complementary variant summand.
Let $p$ be the smallest prime divisor of $n$, and put
\[
 c=n(n-n/p)(g-1).
\]
Proposition~1.4 of \cite{dCMS2022}, which applies to arbitrary rank,
gives
\[
 H^i_{\mathrm{var}}(X)=P_{i-c}H^i_{\mathrm{var}}(X).
\]
Here $c\ge n^2/2\ge8$, so \eqref{norm:perverse-range} makes the
variant summand vanish in degrees $0$, $1$, and $2$.

The invariant summand is the rational cohomology of the corresponding
$\mathrm{PGL}_n$ moduli space.  Markman's generation theorem \cite[Theorem~7]{Markman2002}, in the
projective-bundle formulation of \cite[Theorem~2.1(i)]{MS2024}, states
that this algebra is generated by the K\"unneth components of the
projective Chern characters with indices $j\ge2$.  These components
have degrees $2j$, $2j-1$, and $2j-2$ on the moduli space.  There are
therefore no degree-one generators and only one possible degree-two
generator, namely a nonzero normalization of the expression defining
$e$.  This proves $H^0(X)=\mathbb Q$, $H^1(X)=0$, and
$H^2(X)=\mathbb Q e$, with nonvanishing still to be checked.

Choose an ample line bundle on $X$.  By Lemma~\ref{norm:dimensions},
a Hitchin fiber is projective of positive dimension, and hence contains
an integral projective curve.  The ample bundle has positive degree on
that curve, so its first Chern class is nonzero in $H^2(X)$.  The
one-dimensional upper bound just proved now gives $e\ne0$ and the
final assertion.
\end{proof}

\begin{proposition}\label{norm:relative-lefschetz}
Cup product by $e$ sends $P_kH^m(X)$ into $P_{k+2}H^{m+2}(X)$.
For every $i\ge0$ and every $m$, it induces an isomorphism
\begin{equation}\label{norm:relative-hl}
 e^i:\operatorname{Gr}^P_{R-i}H^m(X)
 \xrightarrow{\ \sim\ }
 \operatorname{Gr}^P_{R+i}H^{m+2i}(X).
\end{equation}
\end{proposition}

\begin{proof}
The filtration shift follows from \eqref{norm:shift-rule} applied to
the cup-product morphism.  Choose an ample class on $X$, which is
also relatively ample for the proper morphism $h$, making $h$
projective.  The decomposition theorem and relative hard Lefschetz
\cite[Theorems~6.2.5 and~6.2.10]{BBD1982} give \eqref{norm:relative-hl} for that
class: the perverse degrees of $Rh_*\mathbb Q_X[2R]$ have center zero,
and the shift to $K$ moves the center to $R$.  The decomposition theorem
also makes the perverse cohomology spectral sequence degenerate, so
these sheaf-theoretic isomorphisms give the displayed isomorphisms on
associated-graded cohomology.  Lemma~\ref{norm:low-cohomology} replaces
the ample class by a nonzero scalar multiple of $e$, without changing
invertibility.
\end{proof}

\subsection{A common lowering operator determines the filtration}

For a finite graded vector space $V=\bigoplus_k V_k$ and a degree-two
operator $E$, say that $E$ is \emph{Lefschetz with center $r$} if
\[
 E^i:V_{r-i}\xrightarrow{\sim}V_{r+i}
 \quad\text{for every }i\ge0.
\]
Here $r$ is an integer, and grades outside the finite support are zero.
Proposition~\ref{norm:relative-lefschetz} gives this property for the
total perverse associated graded, with center $R$.  The next statement
explains the role of the lowering operator in the proof.
The use of an $\mathfrak{sl}_2$-operator to determine the perverse
filtration is also central to \cite[Propositions~7.9 and~8.2--8.3]{HMMS2025}.
We prove the precise two-filtration criterion needed here, including
the passage from associated gradeds to the filtrations themselves.

\begin{proposition}[Comparison by a common lowering operator]
\label{alg:comparison}
Let $H$ be a finite-dimensional vector space over a field of
characteristic zero, with finite, exhaustive, separated increasing
integer-indexed filtrations $F_\bullet$ and $G_\bullet$.
Suppose that $e,f\in\operatorname{End}(H)$ satisfy, for
$A_\bullet=F_\bullet$ and $A_\bullet=G_\bullet$,
\[
 e(A_k)\subseteq A_{k+2},\qquad f(A_k)\subseteq A_{k-2}.
\]
Assume that the induced raiser on each associated graded is Lefschetz
with the same center $r$, and that
\begin{equation}\label{alg:double-commutator}
 [[f,e],e]=-2e,
 \qquad [u,v]=uv-vu.
\end{equation}
Then $F_k=G_k$ for every integer $k$.
If $H$ has a further grading respected by both filtrations, and $e,f$
have degrees $2,-2$ for that grading, the equality holds degree by
degree.
\end{proposition}

\begin{proof}
Fix either associated graded, and denote its raiser by $E$.
The Lefschetz decomposition expresses it as a sum of primitive strings
\[
 v,Ev,\ldots,E^\ell v,\qquad \deg v=r-\ell.
\]
On each string set
\[
 F_0(E^jv)=j(\ell-j+1)E^{j-1}v.
\]
If $H_0$ acts by $k-r$ on grade $k$, direct calculation gives
\[
 [E,F_0]=H_0,\qquad [H_0,E]=2E,\qquad [H_0,F_0]=-2F_0.
\]
Thus $(E,H_0,F_0)$ is an $\mathfrak{sl}_2$-triple, and
$[[F_0,E],E]=-2E$.

We claim that a degree $-2$ endomorphism satisfying this last identity
must equal $F_0$.  On the endomorphism space, the adjoint
$\mathfrak{sl}_2$-action identifies grading degree $-2$ with
$H_0$-weight $-2$.  Decompose this finite-dimensional representation
into irreducibles.  In any irreducible summand containing weight $-2$,
two raising steps map that weight space nontrivially to weight $2$.
Consequently
\[
 (\operatorname{ad}E)^2:
 \operatorname{End}(\operatorname{Gr}^A H)_{-2}
 \longrightarrow
 \operatorname{End}(\operatorname{Gr}^A H)_{2}
\]
is injective.  Applying this to $\operatorname{Gr}^A(f)-F_0$ and
\eqref{alg:double-commutator} proves the claim.

It follows that the single operator $h_0=[e,f]$ on $H$ preserves both
filtrations and acts on their grade $k$ by $k-r$.  Choose integers
$a\le b$ containing all their nonzero grades.  For either filtration,
\[
 (h_0-(k-r))A_k\subseteq A_{k-1}.
\]
Applying these factors successively from $k=b$ down to $k=a$ gives
\[
 \prod_{k=a}^b(h_0-(k-r))=0.
\]
The polynomial has distinct roots, so $h_0$ is semisimple.
Every $h_0$-stable subspace is therefore the direct sum of its
intersections with the eigenspaces.
If a nonzero eigenvector first belongs to $A_j$, its image in
$\operatorname{Gr}^A_jH$ shows that its eigenvalue is $j-r$.
Since each $A_k$ is $h_0$-stable, this proves
\begin{equation}\label{alg:eigenspace-filtration}
 A_k=\bigoplus_{j\le k}\ker\bigl(h_0-(j-r)\bigr).
\end{equation}
The right side is independent of the choice of filtration, proving
$F_k=G_k$.  In the additionally graded setting, $h_0$ has degree zero,
so its eigenspaces and \eqref{alg:eigenspace-filtration} decompose
degree by degree.
\end{proof}

We also record a cancellation fact for the Betti Lefschetz argument.

\begin{lemma}\label{alg:tensor}
Let $V,W$ be nonzero finite integer-graded vector spaces over a field
of characteristic zero, and let $E_V,E_W$ have degree two.
If $E_V\otimes1+1\otimes E_W$ is Lefschetz with center $r$ on
$V\otimes W$, then $E_V$ and $E_W$ are Lefschetz with centers
$r_V,r_W$, respectively, where $r_V+r_W=r$.
\end{lemma}

\begin{proof}
Choose homogeneous Jordan-string decompositions for the two operators,
which are nilpotent because the gradings have finite support.
The tensor product of any pair of strings is a graded invariant direct
summand, with invariant complement, for the sum operator.  Each
Lefschetz isomorphism is block diagonal with respect to these summands,
so it is an isomorphism on each of them.

For a string whose lowest and highest grades are $p$ and $q$, call
$(p+q)/2$ its center.  The average of the lowest and highest grades of
a tensor product of two strings is the sum of their centers.  Its
Lefschetz symmetry forces that sum to equal $r$, for every pair of
strings.  Fixing one string in either factor shows that all strings
in the other factor have the same center.  A Jordan string is
Lefschetz about its own center, which proves the conclusion.
\end{proof}

We will apply Proposition~\ref{alg:comparison} to total cohomology,
with $F_k=P_k$ and $G_k=W_{2k}$ under nonabelian Hodge theory.
Section~\ref{sec:betti} proves that $e$ is also Lefschetz for the
half-weight grading, with center $R$, and that all odd weight-graded
pieces vanish.  The remaining geometric construction supplies one
operator $f:H^m(X)\to H^{m-2}(X)$ satisfying both lowering bounds
and \eqref{alg:double-commutator}.  The proposition then identifies
the two filtrations in every cohomological degree; the separate
even-weight assertion gives $W_{2k}=W_{2k+1}$.
